\chapter{Cloud Against On-Premises}\label{ch:pl:cloud-against-on-premises}

On 4 November 2021 CME Group and Google Cloud announced a ten-year partnership: CME would move its technology to Google Cloud, starting the
next year with data and clearing services and eventually all of its markets. The same year most trading firms kept every order path in their
own racks, as Book~14 describes. Both decisions can be right, because they are about different workloads. This chapter prices the one where the
question is most often asked and least often computed --- the research cluster --- with every cloud price taken from a dated, published list,
every assumption about owned hardware stated, and chapter~\ref{ch:pl:distributed-compute-and-schedulers}'s simulator run on interruptible
capacity.

\section{What moves to the cloud and what cannot}

A public cloud (Book~14, chapter~16) rents machines by the hour in regions and availability zones the firm does not control. Four questions
sort a workload. How sensitive is it to latency and jitter --- the order path, Book~14 shows, is still measured in the firm's own racks next to
the matching engine. How variable is its demand --- a sweep that needs thousands of cores for a weekend and none on Monday is the cloud's
case. Where is its data --- the next section's question. And who must be able to inspect it --- the last section's.

\begin{definition}[Infrastructure as code]\label{def:pl:cloud-against-on-premises:iac}
\emph{Infrastructure as code}\index{infrastructure as code} describes servers, networks, storage and their permissions in version-controlled
files from which a tool creates, changes and destroys them, so that an environment can be rebuilt, reviewed and compared like software.
\end{definition}

\omterm{def:pl:cloud-against-on-premises:iac}{Infrastructure as code} is what makes elastic capacity usable: a cluster that exists only while a sweep runs must be created and destroyed by a
program, identically every time, and an owned cluster described the same way can be rebuilt after a failure. It applies to both sides of this
chapter's comparison.

\section{The economics of compute: owned, reserved, on demand, interruptible}

\begin{definition}[On-demand, reserved and spot capacity]\label{def:pl:cloud-against-on-premises:capacity}
\emph{On-demand capacity}\index{on-demand capacity} is rented by the second or hour with no commitment, at the provider's list price.
\emph{Reserved capacity}\index{reserved capacity} is committed for one or three years at a lower hourly price, paid whether used or not.
\emph{Spot capacity}\index{spot capacity} is the provider's unused capacity rented at a lower, changing price, which the provider can reclaim
at short notice, interrupting whatever runs on it.
\end{definition}

The chapter prices one cloud instance of 96 hardware threads in one US region from the provider's public price list of 25 September 2026: 4.284
dollars an hour on demand, 2.834 reserved for one year and 1.886 for three years without upfront payment, and 1.737 on the spot market on 28
September. Per thread-hour that is 4.46, 2.95, 1.96 and 1.81 cents. Owned hardware has no such list. The model builds a two-socket server of
512 threads from the list price of its CPU (10\,931 dollars each), assumes the CPUs are half of the server's price, spreads it over four years,
powers it at the US average commercial price of electricity (14.53 cents a kilowatt-hour in July 2026) with a 1.3~kW draw and a building
overhead of 1.4, and adds \$300 a month of rack space and a two-hundredth of an engineer's \$250\,000 a year. Those last figures are
assumptions, not facts, and the results below are shown for server prices up to three times the assumption.

\omcode{code/firm/cloudcost/firm_cloudcost.py}{45}{61}{The owned server's monthly cost: purchase spread over its life, power with the
building's overhead, space and staff; and its cost per thread-hour when fully used.}\label{lst:pl:cloud-against-on-premises:owned}

The owned server then costs \$1\,508 a month, 0.40 cents per thread-hour when fully used. The comparison turns on utilisation: capacity that is
paid for every hour --- owned or reserved --- costs its hourly price divided by the share of hours it is used, while on-demand and \omterm{def:pl:cloud-against-on-premises:capacity}{spot capacity}
cost their price only when used (\cref{fig:pl:cloud-against-on-premises:curves}).

\begin{omfigure}
\begin{tikzpicture}
\begin{semilogyaxis}[width=11.5cm,height=6cm,xlabel={utilisation of the capacity paid for},ylabel={cents per used thread-hour (log scale)},
  grid=major,grid style={gray!20},tick label style={font=\small},label style={font=\small},table/col sep=comma,
  xmin=0.05,xmax=1,ymin=0.3,ymax=40,xtick={0.1,0.2,0.4,0.6,0.8,1},
  xticklabel={\pgfmathparse{100*\tick}\pgfmathprintnumber[fixed,precision=0]{\pgfmathresult}\%},
  yticklabel={\pgfmathparse{exp(\tick)}\pgfmathprintnumber[fixed,precision=1]{\pgfmathresult}}]
\addplot[omThm,thick] table[x=u,y expr=100*\thisrow{owned}]{figdata/platforms/26-cloud-against-on-premises/curves.csv};
\addplot[omDef,thick,dashed] table[x=u,y expr=100*\thisrow{reserved_3y}]{figdata/platforms/26-cloud-against-on-premises/curves.csv};
\addplot[black,thick] table[x=u,y expr=100*\thisrow{on_demand}]{figdata/platforms/26-cloud-against-on-premises/curves.csv};
\addplot[gray,thick,dashdotted] table[x=u,y expr=100*\thisrow{spot}]{figdata/platforms/26-cloud-against-on-premises/curves.csv};
\draw[omThm,dotted] (axis cs:0.090,0.3) -- (axis cs:0.090,4.46) node[pos=0.25,right,font=\scriptsize]{9\%};
\draw[omThm,dotted] (axis cs:0.223,0.3) -- (axis cs:0.223,1.81) node[pos=0.1,right,font=\scriptsize]{22\%};
\node[omThm,font=\scriptsize,anchor=south west] at (axis cs:0.6,0.68) {owned};
\node[omDef,font=\scriptsize,anchor=south west] at (axis cs:0.75,2.8) {reserved 3 years};
\node[font=\scriptsize,anchor=south west] at (axis cs:0.6,4.6) {on demand};
\node[gray,font=\scriptsize,anchor=north west] at (axis cs:0.6,1.75) {spot, no preemption loss};
\end{semilogyaxis}
\end{tikzpicture}

\omcaption{Cost per thread-hour actually used against the utilisation of the capacity paid for, at the cited prices and the model's owned
server. Owned capacity is cheaper than on-demand above 9\% utilisation and cheaper than spot above 22\%; three-year \omterm{def:pl:cloud-against-on-premises:capacity}{reserved capacity} beats
on-demand above 44\% but never beats owning. Data: \texttt{fig\_cloudcost.py}.}\label{fig:pl:cloud-against-on-premises:curves}
\end{omfigure}

\begin{table}[htbp]
\centering\small
\begin{tabular}{lrrrr}
\toprule
server price & per month & cents per thread-hour & break-even: on demand & spot\\
\midrule
as assumed (\$43\,724) & \$1\,508 & 0.40 & 9.0\% & 22.3\%\\
twice & \$2\,419 & 0.65 & 14.5\% & 35.8\%\\
three times & \$3\,330 & 0.89 & 20.0\% & 49.2\%\\
\bottomrule
\end{tabular}
\caption{The utilisation above which the owned server costs less than renting the same threads, for three server prices. Even at three times
the assumed price, owning wins for any capacity used more than a fifth of the time against on-demand, and half the time against spot.}
\label{tab:pl:cloud-against-on-premises:break-even}
\end{table}

A thread is not a unit of work: the owned CPU and the cloud instance's CPU run a given backtest at different speeds, and a firm should compare
them with chapter~\ref{ch:pl:accelerators-for-quantitative-work}'s measurements of its own workload before trusting any break-even.
A factor of two in speed divides or multiplies the owned server's whole cost per unit of work by two, and the break-even with it: 4.5\% or
18\% against on demand.

\section{Buying for a year of demand}

A research cluster's demand is not a utilisation but a profile. The model's year has an interactive base of 1\,500 threads, 1\,000 more in
weekday working hours, a nightly batch of 4\,000 more on weekdays, a weekend sweep of 12\,000 more for eight hours every Saturday, and twenty
unplanned bursts of 8\,000 for four hours: a mean of 3\,364 threads against a peak of 13\,500. Sorting the hours by demand gives the load-duration
curve, and each level of capacity is used in the share of hours the demand reaches it.

\omcode{code/firm/cloudcost/firm_cloudcost.py}{88}{105}{The capacity planner: the $k$-th unit of capacity is needed in the share $f(k)$ of
hours; it is owned, reserved or bought by the hour, whichever costs least for that share.}\label{lst:pl:cloud-against-on-premises:plan}

\begin{omfigure}
\begin{tikzpicture}
\begin{axis}[width=11.5cm,height=5.6cm,xlabel={hours of the year, sorted by demand},ylabel={threads},
  grid=major,grid style={gray!20},tick label style={font=\small},label style={font=\small},table/col sep=comma,
  xmin=0,xmax=8760,ymin=0,ymax=15000,scaled y ticks=false,yticklabel style={/pgf/number format/1000 sep={\,}},
  xticklabel style={/pgf/number format/1000 sep={\,}}]
\addplot[omDef,thick,const plot] table[x=hours,y=threads]{figdata/platforms/26-cloud-against-on-premises/duration.csv};
\fill[omThm!15] (axis cs:0,0) rectangle (axis cs:8760,5500);
\addplot[omDef,thick,const plot] table[x=hours,y=threads]{figdata/platforms/26-cloud-against-on-premises/duration.csv};
\draw[omThm,dashed] (axis cs:788,0) -- (axis cs:788,15000) node[pos=0.93,right,font=\scriptsize]{9\% of hours};
\node[omThm,font=\scriptsize] at (axis cs:5000,3000) {owned: 5\,500 threads (about 11 servers)};
\node[font=\scriptsize,anchor=south west] at (axis cs:900,9500) {bought by the hour: 496 hours above 5\,500};
\end{axis}
\end{tikzpicture}

\omcaption{The model year's load-duration curve. Every level of demand reached in more than 9\% of hours --- the break-even against
on-demand --- is cheaper owned: here up to the 5\,500 threads of the weekday nightly batch. The weekend sweeps and unplanned bursts above it
occupy 496 hours and are bought by the hour. Data: \texttt{fig\_cloudcost.py}.}\label{fig:pl:cloud-against-on-premises:duration}
\end{omfigure}

\begin{omfigure}
\begin{tikzpicture}
\begin{axis}[width=11.5cm,height=5.4cm,xbar stacked,bar width=11pt,xmin=0,xmax=1500,y dir=reverse,
  ytick=data,yticklabels={on demand only,reserved and on demand,owned and on demand,owned and spot bursts},
  xlabel={annual cost (thousand dollars)},grid=major,grid style={gray!20},tick label style={font=\small},label style={font=\small},
  table/col sep=comma,legend style={font=\small,at={(0.5,1.03)},anchor=south,legend columns=3,draw=none,
  /tikz/every even column/.append style={column sep=3mm}}]
\addplot[fill=omThm!60,draw=omThm] table[y expr=\coordindex,x expr=\thisrow{owned}/1000]{figdata/platforms/26-cloud-against-on-premises/plans.csv};
\addplot[fill=omDef!40,draw=omDef] table[y expr=\coordindex,x expr=\thisrow{reserved}/1000]{figdata/platforms/26-cloud-against-on-premises/plans.csv};
\addplot[fill=gray!40,draw=gray] table[y expr=\coordindex,x expr=\thisrow{metered}/1000]{figdata/platforms/26-cloud-against-on-premises/plans.csv};
\legend{owned,reserved (3 years),by the hour}
\node[font=\scriptsize,anchor=west] at (axis cs:1320,0) {1\,315};
\node[font=\scriptsize,anchor=west] at (axis cs:945,1) {940};
\node[font=\scriptsize,anchor=west] at (axis cs:366,2) {361};
\node[font=\scriptsize,anchor=west] at (axis cs:287,3) {282};
\end{axis}
\end{tikzpicture}

\omcaption{The model year's demand bought four ways. On demand alone costs \$1.32 million; reserving the 2\,500 threads used most of the time
brings it to \$0.94 million; owning 5\,500 threads and buying the rest on demand, to \$0.36 million; running 80\% of the bursts on spot, to
\$0.28 million. Data: \texttt{fig\_cloudcost.py}.}\label{fig:pl:cloud-against-on-premises:plans}
\end{omfigure}

\begin{definition}[Cloud bursting]\label{def:pl:cloud-against-on-premises:bursting}
\emph{Cloud bursting}\index{cloud bursting} runs a workload on the firm's own capacity up to its size and sends the excess to rented cloud
capacity, so that the owned cluster is sized for the frequent load and the cloud absorbs the rare peaks.
\end{definition}

The optimum in \cref{fig:pl:cloud-against-on-premises:plans} is a hybrid, and it is one because of the shape of the demand, not because of any
single price: the base and the nightly batch are used more than 9\% of the time and are cheaper owned; the weekend sweeps and bursts, 496 hours
of the year, are cheaper rented. A firm whose research runs flat out every night owns; a firm whose research is a few enormous sweeps a year
rents; most are in between.

\section{Spot capacity and the cost of interruption}

\omterm{def:pl:cloud-against-on-premises:capacity}{Spot capacity} is 59\% cheaper than on demand for this instance, and the provider's own advisor places it in the highest band of interruption
frequency, above 20\% of instances reclaimed in the trailing month --- at least 0.03\% an hour. Chapter~\ref{ch:pl:distributed-compute-and-schedulers}'s
sweep of 10\,000 backtests, forty of them seven and a half hours long, runs on 24 instances, each billed until its last task ends.

\omcode{code/firm/cloudcost/firm_cloudcost.py}{109}{124}{Checkpointing as a task graph: each backtest cut into chained chunks, so that an
interruption loses at most one chunk.}\label{lst:pl:cloud-against-on-premises:ckpt}

\begin{omfigure}
\begin{tikzpicture}
\begin{semilogxaxis}[width=11.5cm,height=5.6cm,xlabel={interruptions per instance-hour (log scale)},ylabel={cost of the sweep (\$)},
  grid=major,grid style={gray!20},tick label style={font=\small},label style={font=\small},table/col sep=comma,
  xmin=0.0002,xmax=0.3,ymin=0,ymax=220]
\addplot[omThm,thick,mark=*] table[x=rate_per_hour,y=cost_none]{figdata/platforms/26-cloud-against-on-premises/burst.csv};
\addplot[omDef,thick,mark=square*] table[x=rate_per_hour,y=cost_ckpt]{figdata/platforms/26-cloud-against-on-premises/burst.csv};
\draw[black,dashed] (axis cs:0.0002,202.79) -- (axis cs:0.3,202.79) node[pos=0.02,below right,font=\scriptsize]{on demand: \$203, 7.5 hours};
\draw[gray,dotted] (axis cs:0.000306,0) -- (axis cs:0.000306,220) node[pos=0.55,right,font=\scriptsize,align=left]{advisor's\\highest band};
\node[omThm,font=\scriptsize,anchor=south east] at (axis cs:0.2,150) {spot, restart from scratch};
\node[omThm,font=\scriptsize,anchor=east] at (axis cs:0.17,184) {59 hours};
\node[omDef,font=\scriptsize,anchor=north west] at (axis cs:0.0005,72) {spot, 15-minute checkpoints: \$80--82, under 9 hours};
\end{semilogxaxis}
\end{tikzpicture}

\omcaption{Chapter~13's sweep on 24 spot instances against the interruption rate (log scale; each instance billed until its last task ends).
Without checkpoints, interruptions restart the long backtests from the beginning: at the advisor's band the sweep costs \$94 and takes 14.4
hours instead of 7.5; at 0.2 an hour, \$185 and 59 hours. With 15-minute checkpoints it costs \$80--82 and finishes in under 9 hours at every
rate. Data: \texttt{fig\_cloudcost.py}.}\label{fig:pl:cloud-against-on-premises:burst}
\end{omfigure}

On demand, the sweep costs \$203 and takes 7.5 hours, the length of its longest backtest. On spot without interruptions it costs \$82. What
interruption costs is mostly time, and it falls on the long tasks (\cref{fig:pl:cloud-against-on-premises:burst}): a 7.5-hour backtest
interrupted after six hours starts again, so a single interruption of one of the forty long tasks nearly doubles the sweep's length. Cutting
each backtest into 15-minute chunks that write a checkpoint --- the training checkpoint of Book~12, chapter~23, applied to backtests --- bounds the
loss at one chunk, and the sweep costs the same at every interruption rate the chart shows. \omterm{def:pl:cloud-against-on-premises:capacity}{Spot capacity} is for work that can be interrupted,
and checkpointing is what makes work interruptible.

\section{Data gravity and data-transfer charges}

\begin{definition}[Data gravity]\label{def:pl:cloud-against-on-premises:gravity}
\emph{Data gravity}\index{data gravity} is the tendency of computation to move to where large data already is, because moving the data costs
more in time and data-transfer charges than moving the computation.
\end{definition}

The provider's data-transfer charge (Book~14, chapter~16) is asymmetric: data coming into the region is free, data leaving it to the internet
costs 9 cents a gigabyte for the first 10 terabytes of a month, falling in tiers to 5 cents above 150 terabytes. Keeping a 500-terabyte \omterm{def:pl:a-tick-store-on-open-formats:store}{tick
store} in the cloud's standard object storage (chapter~\ref{ch:pl:capturing-and-storing-tick-data}'s price) costs \$11\,776 a month; taking a
full copy out once costs \$29\,491, about twenty months of an owned server. So the \omterm{def:pl:a-tick-store-on-open-formats:store}{tick store}'s location decides where research runs: once
it is in the cloud, research follows it, and results --- a few gigabytes a sweep --- are what should cross the boundary. A hybrid design keeps
one authoritative copy and sends the computation to it, or keeps two copies and pays once to fill the second.

\section{Security, control and regulators}

\begin{definition}[Shared responsibility model]\label{def:pl:cloud-against-on-premises:shared}
The \emph{shared responsibility model}\index{shared responsibility model} divides security between a cloud provider, responsible for the
infrastructure that runs its services, and the customer, responsible for what it builds and configures on them --- data, identities, network
rules, encryption --- to an extent that depends on the services it chooses.
\end{definition}

Most cloud incidents at financial firms are on the customer's side of that line: a storage bucket left public, a key committed to a repository,
a permission granted too widely --- the subjects of chapter~\ref{ch:pl:security}. Regulators add a second line. In the European Union the EBA's
guidelines on outsourcing, which apply since 30 September 2019 and absorbed its earlier recommendation on cloud outsourcing, require a
documented exit strategy for every critical or important function outsourced: a plan, tested, for moving the function elsewhere if the provider
fails or the relationship ends. For a research cluster that is easy --- the owned half of a hybrid is the exit. For an order path or a clearing
system, it is a second full implementation, which is part of why a move like CME's is a ten-year programme. Concentration --- many firms relying
on the same few providers --- is the operational risk (Book~6, chapter~28) that supervisors watch most.

\begin{dated}{2026-09}{Cloud prices and the rules that apply}\label{dat:pl:cloud-against-on-premises:prices}
AWS's public price list of 25 September 2026 gives, for a c7i.24xlarge instance (96 vCPUs) in US East (N. Virginia) running Linux, 4.284 dollars an
hour on demand, 2.83387 reserved for one year and 1.88608 for three years without upfront payment; its spot price feed gave 1.7372 on 28
September, and its Spot Instance Advisor put the instance's interruption frequency in the highest band (above 20\% in the trailing month) with
65\% savings over on demand. Data transfer out of the region to the internet costs 0.09 dollars per gigabyte for the first 10~TB a month, 0.085
for the next 40, 0.07 for the next 100 and 0.05 above 150~TB; inbound is free. AMD lists its 128-core EPYC 9755 at 10\,931 dollars (1\,000-unit
price); the EIA reports an average US commercial electricity price of 14.53 cents per kWh for July 2026. CME Group and Google Cloud announced
their ten-year partnership, with a one-billion-dollar equity investment by Google, on 4 November 2021. The EBA's outsourcing guidelines
(EBA/GL/2019/02) apply from 30 September 2019.
\end{dated}

\section{Tutorial: own, rent or both}\label{tut:pl:cloud-against-on-premises:tut}

\begin{tutorial}
\textbf{Goal.} Price a research cluster's year every way, find the break-even utilisation, and run a sweep on \omterm{def:pl:cloud-against-on-premises:capacity}{spot capacity}.
\textbf{End state:} \cref{fig:pl:cloud-against-on-premises:curves,fig:pl:cloud-against-on-premises:plans,fig:pl:cloud-against-on-premises:burst}.
\begin{tutsteps}
\item \textbf{Prices}: \texttt{pl\_cloudcost.PRICES} and \texttt{SERVER}; \texttt{per\_thread()} and \texttt{break\_evens(k)} for
$k=1,2,3$.
\item \textbf{Demand}: \texttt{demand()}; plot its load-duration curve.
\item \textbf{Plans}: \texttt{plans()}, the four ways of buying the year.
\item \textbf{Burst}: \texttt{burst(rate, price, checkpoint)} over the interruption rates, with and without checkpoints.
\item \textbf{Data}: \texttt{firm\_cloudcost.transfer\_out} for the \omterm{def:pl:a-tick-store-on-open-formats:store}{tick store} and for a sweep's results.
\end{tutsteps}
\textbf{What to change next.} Replace the thread as the unit by measured backtests per hour on each machine; let \omterm{def:pl:cloud-against-on-premises:capacity}{reserved capacity} cover the
nightly batch when the firm cannot buy servers in time.
\end{tutorial}

\section{Build: the cloud cost model}\label{bld:pl:cloud-against-on-premises:cloudcost}

\begin{build}
\bfield{Purpose} Price any compute plan from dated, cited lists and stated assumptions, and choose the mix of owned, reserved, on-demand and
\omterm{def:pl:cloud-against-on-premises:capacity}{spot capacity} for a demand profile.
\bfield{Interface} \texttt{PriceItem}, \texttt{OwnedServer}, \texttt{cost\_per\_used\_hour}, \texttt{break\_even}, \texttt{plan\_capacity},
\texttt{checkpointed}, \texttt{billed\_node\_hours}, \texttt{transfer\_out}.
\bfield{Rules} Every price carries its source and date; every assumption is named and shown with a sensitivity; plans are compared on the
same demand; interruptible capacity only for checkpointed work; data-transfer charges counted both ways.
\bfield{Acceptance tests} \texttt{code/firm/cloudcost/tests/}: the owned server's monthly and hourly cost; fixed and metered costs and the
break-even; the planner on a two-level demand, with reserved and \omterm{def:pl:cloud-against-on-premises:capacity}{spot capacity}; checkpoint chunks, their chaining and the billing; the tiers of
the transfer charge.
\bfield{Stretch} Savings plans and committed-use discounts; interruption rates from measured history; \omterm{def:pl:capturing-and-storing-tick-data:tier}{storage tiers} and retrieval charges in
the data-gravity comparison.
\end{build}

\begin{omsources}
\item CME Group and Google Cloud, press release of 4 November 2021.
\item AWS public price lists (EC2 on demand and reserved, data transfer), spot price feed and Spot Instance Advisor; AWS, \emph{Shared
Responsibility Model}.
\item AMD, EPYC 9755 product page; U.S. Energy Information Administration, \emph{Electric Power Monthly}, Table 5.6.A.
\item European Banking Authority, \emph{Guidelines on outsourcing arrangements}, EBA/GL/2019/02.
\item One Quant Book~14, chapter~16 (public clouds, zones, data-transfer charges); Book~16, chapter~20 (total cost of ownership).
\end{omsources}

\section{Exercises}

\begin{exercise}[$\star$]\label{exo:pl:cloud-against-on-premises:1}
What does one thread-hour cost on demand, reserved for three years, and on spot, at the cited prices?
\end{exercise}

\begin{exercise}[$\star$]\label{exo:pl:cloud-against-on-premises:2}
Above what utilisation is three-year \omterm{def:pl:cloud-against-on-premises:capacity}{reserved capacity} cheaper than on-demand?
\end{exercise}

\begin{exercise}[$\star$]\label{exo:pl:cloud-against-on-premises:3}
What does it cost to take a 500-terabyte \omterm{def:pl:a-tick-store-on-open-formats:store}{tick store} out of the region once?
\end{exercise}

\begin{exercise}[$\star\star$]\label{exo:pl:cloud-against-on-premises:4}
Why does the planner never choose \omterm{def:pl:cloud-against-on-premises:capacity}{reserved capacity} when owning is allowed?
\end{exercise}

\begin{exercise}[$\star\star$]\label{exo:pl:cloud-against-on-premises:5}
Why does a single interruption of a long backtest nearly double the sweep's length?
\end{exercise}

\begin{exercise}[$\star\star$]\label{exo:pl:cloud-against-on-premises:6}
The firm's CPUs run a backtest twice as fast as the cloud instance's. What happens to the break-even?
\end{exercise}

\begin{exercise}[$\star\star\star$]\label{exo:pl:cloud-against-on-premises:7}
\emph{Coding.} Rerun \texttt{plans()} on \texttt{demand(spread\_sweep=True)}, where the weekend sweep's thread-hours run inside the five
weeknight \omterm{def:pl:the-risk-grid:grid}{batch windows} instead. How much capacity does the planner own, and what does the year cost?
\end{exercise}

\begin{exercise}[$\star\star\star$]\label{exo:pl:cloud-against-on-premises:8}
\emph{Find the flaw.} \textquotedbl{}Spot is 59\% cheaper, so we will run the whole research cluster on spot.\textquotedbl{}
\end{exercise}

\section{Problem: Own, Rent or Both}

\begin{problem}[{Weekend problem --- own, rent or both}]\label{pb:pl:cloud-against-on-premises:1}
The chapter's prices, owned server, demand profile and sweep.

\medskip\noindent\textbf{Part I --- Prices.}
\begin{enumerate}
\item What are the four cloud prices per thread-hour, and where do they come from?
\item What is cited and what is assumed in the owned server's cost?
\item What does the owned server cost a month and per thread-hour?
\item How does utilisation enter the comparison?
\item What are the break-evens against on-demand and spot, and how do they move with the server's price?
\end{enumerate}

\medskip\noindent\textbf{Part II --- The year.}
\begin{enumerate}[resume]
\item What is the model's demand profile, its mean and its peak?
\item How does the planner decide each unit of capacity?
\item How much does it own, and how many hours are bought by the hour?
\item What does each of the four plans cost?
\item Why is the optimum a hybrid?
\end{enumerate}

\medskip\noindent\textbf{Part III --- Spot and data.}
\begin{enumerate}[resume]
\item What does the provider say about this instance's interruptions?
\item What do the sweep's cost and length become on spot without checkpoints?
\item What do checkpoints change?
\item What does it cost to store the \omterm{def:pl:a-tick-store-on-open-formats:store}{tick store} in the cloud, and to take it out?
\item What does a regulator require before a critical function moves to a cloud?
\end{enumerate}

\medskip\noindent\textbf{Part IV --- The verdict.}
\begin{enumerate}[resume]
\item State the \emph{named result}: the utilisation above which owning a research cluster costs less than renting it at cited prices, and
the cost of the chapter~13 sweep on \omterm{def:pl:cloud-against-on-premises:capacity}{spot capacity} with preemption against on demand.
\item Which of the model's assumptions would you check first, and how?
\item Where should the \omterm{def:pl:a-tick-store-on-open-formats:store}{tick store} live?
\item What would make you move the whole cluster to the cloud?
\item In one sentence: what decides between owning and renting compute?
\end{enumerate}
\end{problem}

\section{Interview questions}

\begin{interviewq}[$\star$ \iqroles{developer}]\label{iq:pl:cloud-against-on-premises:1}
What are on-demand, reserved and spot instances, and when would you use each?
\end{interviewq}

\begin{interviewq}[$\star\star$ \iqroles{developer}]\label{iq:pl:cloud-against-on-premises:2}
How would you decide whether to buy a research cluster or rent one?
\end{interviewq}

\begin{interviewq}[$\star\star$ \iqroles{developer}]\label{iq:pl:cloud-against-on-premises:3}
How do you run a large backtest sweep on capacity that can be taken away at any moment?
\end{interviewq}

\begin{interviewq}[$\star\star$ \iqroles{developer}]\label{iq:pl:cloud-against-on-premises:4}
Your \omterm{def:pl:a-tick-store-on-open-formats:store}{tick store} is 500 terabytes on premises. A team wants to run research in the cloud. What do you tell them?
\end{interviewq}

\begin{interviewq}[$\star\star\star$ \iqroles{developer}]\label{iq:pl:cloud-against-on-premises:5}
Design a hybrid research platform that uses owned capacity for the base load and the cloud for peaks.
\end{interviewq}
